\documentclass[a6paper,twoside]{article}

% morewrites virtualises \newwrite so a document that combines many glossary
% groups (each costs one \write), biblatex (\bcfout), hyperref (\@outlinefile),
% the toc, and the LOT/LOF can exceed XeLaTeX's 16-register limit. Must be
% loaded before any package that allocates writes; that is why it sits ahead
% of xcolor here.
\usepackage{morewrites}

\usepackage[svgnames,dvipsnames,x11names]{xcolor}
\usepackage{graphicx}
% Babel must be loaded before any package that hooks into its localisation (notably
% the glossaries fragment in ``extra_packages``), otherwise the default acronym
% table title falls back to English regardless of ``language: french`` etc.
\usepackage[english]{babel}
% Hyperref must be loaded before ``glossaries`` (pulled in via ``extra_packages``)
% so glossaries can wrap glossary page numbers in ``\hyperpage`` for clickable
% backreferences. Loading it after glossaries leaves the page numbers as plain
% text without hyperlinks.
\usepackage[colorlinks=true, linkcolor=NavyBlue, urlcolor=NavyBlue, citecolor=NavyBlue]{hyperref}
\usepackage[a6paper]{geometry}
\usepackage{ts-fonts}
\usepackage{float}
\usepackage{seqsplit}
\usepackage{ts-index}

\usepackage{lastpage} % pour le label LastPage
\usepackage{refcount} % pour \getpagerefnumber
\makeatletter
\def\ps@plain{%
  \let\@oddhead\@empty
  \let\@evenhead\@empty
  \def\@oddfoot{%
    \hfil
    % Afficher le numéro seulement si le document a plus d'une page
    \ifnum\getpagerefnumber{LastPage}>1\relax
      \thepage
    \fi
    \hfil}%
  \let\@evenfoot\@oddfoot
}
\makeatother

\title{Index Example}
\author{}\date{}

\begin{document}

\maketitle

\thispagestyle{plain}
This document demonstrates the index generation with TeXSmith. You can nest, format and create multiple index entries for terms.

\section{Granny Smith}\label{granny-smith}

A tart apple variety often used in baking and cooking. Known for its bright green skin and crisp texture.
\index{Granny Smith} \index{apple}

\clearpage

\section{Fuji}\label{fuji}

A sweet and juicy apple variety that originated in Japan. It has a dense flesh and a balanced flavor.

\index{Fuji} \index{apple} \index{juicy@\textit{juicy}}

\clearpage

\section{Honeycrisp}\label{honeycrisp}

A popular apple variety known for its crisp texture and sweet-tart flavor. It has a distinctive red and yellow skin.

\index{Honeycrisp} \index{apple}

\printindex

\end{document}
